\section{Game bijections and isomorphisms}
For games $G=(I,(A_i),(u_i))$ and $H=(J,(B_j),(v_j))$, a game bijection
$g=(\pi,(g_i)_{i\in I})$ consists of a bijection $\pi:I\to J$ and
bijections $g_i:A_i\to B_{\pi(i)}$. Its profile map is
\[
 (gs)_{\pi(i)}=g_i(s_i).
\]
It is an \emph{exact game isomorphism} when
$u_i(s)=v_{\pi(i)}(gs)$ for all $i,s$.
This is Ham's Definition~4.8 \citep{ham2021}; transformations preserving
only preferences are not included in this definition.

If $h=(\rho,(h_j)):H\to K$, define
\[
 hg=(\rho\pi,(h_{\pi(i)}\circ g_i)_{i\in I}),\qquad
 g^{-1}=(\pi^{-1},(g_{\pi^{-1}(j)}^{-1})_{j\in J}).
\]
These maps have the required domains and satisfy $(hg)s=h(gs)$.
Composition is associative because the underlying player and strategy maps
compose associatively. The identity uses identity maps in every component.
For isomorphisms, substituting payoff equalities shows that composition and
inversion preserve the isomorphism condition. This verifies the groupoid
constructions of Ham's Theorems~4.7 and~4.9 \citep{ham2021}.

An automorphism is an isomorphism from a game to itself. Its player component
defines a homomorphism $q:\Aut(G)\to\Sym(I)$. An automorphism fixing each
player may still relabel their strategies. Consequently the player image
$q(\Aut(G))$ and the full automorphism group must be distinguished.
